\chapter{Releases and continuous integration}

\eyebrow{how versions are built and shipped}

\vspace{4pt}
The project uses automated quality checks and multi-platform builds. Two GitHub
Actions workflows implement this: one verifies every change, and one produces a
release for every version tag.

\section{Continuous integration}

The \textbf{CI} workflow runs on every push and pull request to the main branch.
It verifies the three languages of the project in parallel:

\begin{center}
\begin{tabularx}{\linewidth}{@{}l X@{}}
\toprule
\textbf{Job} & \textbf{What it checks} \\
\midrule
Go core & \tele{go vet} and the full test suite (\tele{go test ./...}), including
      the pairing-integrity and Fleiss'~$\kappa$ tests. \\
Wails build & A complete \tele{wails build}: it generates the Go$\leftrightarrow$JS
      bindings, type-checks the front end (\tele{tsc}), builds it (\tele{vite}) and
      compiles the Go core, validating all three layers together, as a release
      does. \\
Python sidecars & Byte-compile (syntax) plus a \tele{ruff} lint of the inference,
      conformation, live and Kinect sidecars. \\
\bottomrule
\end{tabularx}
\end{center}

A green CI badge on the repository means all three passed on the current main
branch.

\section{Releases}

The \textbf{Release} workflow triggers when a version tag matching \tele{v*} is
pushed. It builds the native application on three runners in parallel (macOS,
Windows and Linux). It then bundles the model weights into each build, and
attaches the archives to a GitHub Release with auto-generated notes.

\begin{center}
\begin{tabularx}{\linewidth}{@{}l l@{}}
\toprule
\textbf{Platform} & \textbf{Release archive} \\
\midrule
macOS (universal) & \tele{carcass-macos.tar.gz} \\
Windows (amd64)   & \tele{carcass-windows.zip} \\
Linux (amd64)     & \tele{carcass-linux.tar.gz} \\
\bottomrule
\end{tabularx}
\end{center}

Each archive is self-contained: the trained weights (\tele{fat\_binary.pth},
\tele{direct\_class.pth}, \tele{eg\_regression.pth}) and the ColorNaming lookup
table travel with the binary, placed where the app expects them at runtime.

\section{Publishing a new version}

A release is triggered by a single tagged push:

\begin{lstlisting}
git tag -a v0.2.0 -m "Apex v0.2.0 - <summary>"
git push origin v0.2.0
\end{lstlisting}

The workflow then performs the three platform builds, the weight bundling, and the
published release with downloadable binaries. The version in \tele{wails.json}
should be kept in step with the tag.

\begin{note}[title={Reproducible builds}]
Because the release uses the same \tele{wails build} that the CI runs on every
commit, a green CI indicates that the release will build. The result is a set of
ready-to-run binaries that need no build tools. The model features require a
Python environment with the scientific libraries.
\end{note}
